\documentclass{article}
\usepackage{geometry}
\geometry{margin=1.5cm, vmargin={0pt,1cm}}
\setlength{\topmargin}{-1cm}
\setlength{\headheight}{12.5pt}
\setlength{\paperheight}{29.7cm}
\setlength{\textheight}{25.3cm}

\usepackage[UTF8]{ctex}
\usepackage{amsmath,amssymb,amsthm,amsfonts}
\usepackage{enumerate}
\usepackage{fancyhdr}
\usepackage{xcolor}
\usepackage{hyperref}
\hypersetup{colorlinks=true, linkcolor=blue!60!black, urlcolor=blue!60!black}

\newcommand{\C}{\mathbb{C}}
\newcommand{\R}{\mathbb{R}}
\newcommand{\Z}{\mathbb{Z}}
\newcommand{\N}{\mathbb{N}}
\newcommand{\dd}{\mathrm{d}}
\newcommand{\Res}{\operatorname{Res}}
\newcommand{\Aut}{\operatorname{Aut}}
\renewcommand{\Re}{\operatorname{Re}}
\renewcommand{\Im}{\operatorname{Im}}
\DeclareMathOperator{\Log}{Log}
\DeclareMathOperator{\Arg}{Arg}

\newcommand{\sourceHandout}{\texttt{complex\_analysis/handouts/bookComplexAnalysisHandout-2026-06-22.pdf}}
\newcommand{\keyFile}[1]{\texttt{#1}}

\begin{document}

\pagestyle{fancy}
\fancyhead{}
\lhead{Complex Analysis --- Chapter Summary}
\rhead{\today}

\begin{center}
  {\Large\bfseries 复变函数按章复习整理}\\[6pt]
  {\large 主要内容、重要定理与典型题目分析}
\end{center}
\vspace{6pt}

\noindent
依据最新版整本 handout：\sourceHandout。
本文档按讲义 5 章结构组织，并补充每章常见题型和解题路线。

\tableofcontents
\newpage

%=====================================================================
\section*{说明}
\addcontentsline{toc}{section}{说明}
%=====================================================================

依据：\sourceHandout，即目前 \keyFile{handouts} 目录中日期最新的复变函数整本 handout。

已有类似文件：

\begin{itemize}
\item
  \texttt{complex\_analysis/handouts/key-points.md}：已有一份中文知识点提纲，但主要按旧的分章
  handout 整理，缺少典型题目分析，也未完整覆盖最新版整本 handout 中第 4
  章后半的留数、Mobius 变换和第 5
  章的辐角原理、Rouche、Hurwitz、拓扑次数。
\item
  \texttt{complex\_analysis/review/main.tex} 及同目录若干
  \texttt{ch*.tex}：已有 LaTeX
  版定理/定义复习笔记，但覆盖面偏向前半课程，也不是按最新版整本 handout
  的 5 章结构组织。
\end{itemize}

%=====================================================================
\section{第 1 章：复数与基本函数 (Complex Numbers)}
%=====================================================================

\subsection{主要内容}

本章从数系扩张讲起：自然数、整数、有理数、实数，再到复数域
$\C$。真正需要掌握的是复数域的代数结构、复数列与级数、幂级数收敛半径，以及指数、三角、对数、幂函数等基本复函数。

重点主线如下：

\begin{itemize}
\item
  $\C=\R\times\R$，乘法定义为 \[
  (x+iy)(u+iv)=(xu-yv)+i(xv+yu).
  \]
\item
  共轭、模、实部虚部的关系： \[
  \bar z=x-iy,\quad |z|^2=z\bar z,\quad
  \operatorname{Re} z=\frac{z+\bar z}{2},\quad
  \operatorname{Im} z=\frac{z-\bar z}{2i}.
  \]
\item
  复数列收敛等价于实部、虚部分别收敛。
\item
  函数项级数的正常收敛推出绝对收敛和局部一致收敛。
\item
  幂级数有收敛半径，圆内正常收敛，圆外发散。
\item
  指数函数由幂级数定义，满足加法公式和 $2\pi i$ 周期性。
\item
  复对数是多值的，单值分支必须在割平面上选取。
\item
  复数乘法对应平面相似变换：缩放 $|a|$ 加旋转
  $\arg a$。
\end{itemize}

\subsection{重要定理和结论}

\begin{itemize}
\item
  $\C$ 不是有序域：因为有序域中平方非负，但 $i^2=-1$。
\item
  三角不等式与反三角不等式： \[
  |z+w|\le |z|+|w|,\qquad ||z|-|w||\le |z-w|.
  \]
\item
  Weierstrass M-test：正常收敛推出绝对收敛和局部一致收敛。
\item
  Abel 定理：幂级数存在唯一收敛半径。
\item
  Cauchy-Hadamard 公式： \[
  R=\frac{1}{\limsup_{n\to\infty}|a_n|^{1/n}}.
  \]
\item
  指数函数： \[
  e^{z+w}=e^z e^w,\qquad e^z\ne 0,\qquad e^{z+2\pi i}=e^z.
  \]
\item
  Euler 公式： \[
  e^{iy}=\cos y+i\sin y.
  \]
\item
  主值对数： \[
  \Log z=\log|z|+i\Arg z
  \] 只能在选定割线后的区域中连续解析。
\item
  幂函数一般通过对数分支定义： \[
  z^\alpha=\exp(\alpha \Log z).
  \]
\end{itemize}

\subsection{典型题目分析}

题型 1：求幂级数收敛半径。

常见题面：求 \[
\sum_{n=1}^{\infty}a_n z^n
\] 的收敛半径，其中 $a_n$ 可能包含
$n!$、$n^n$、$\log n$ 或组合式。

分析方法：

\begin{itemize}
\item
  首先尝试比值法： \[
  R=\lim_{n\to\infty}\left|\frac{a_n}{a_{n+1}}\right|
  \] 前提是极限存在。
\item
  比值法不好用时，用 Cauchy-Hadamard： \[
  \frac1R=\limsup |a_n|^{1/n}.
  \]
\item
  遇到阶乘，通常使用 Stirling 公式 \[
  n!\sim \sqrt{2\pi n}\left(\frac ne\right)^n.
  \]
\end{itemize}

典型结论：

\begin{itemize}
\item
  $a_n=(\log n)^2$
  时，$|a_n|^{1/n}\to 1$，所以
  $R=1$。
\item
  $a_n=n!$
  时，$|a_n|^{1/n}\to\infty$，所以
  $R=0$。
\item
  $a_n=1/n!$ 时，$R=\infty$。
\end{itemize}

题型 2：解复指数、复对数和复幂方程。

常见题面：求解 $e^z=w$、$z^n=a$、$i^i$
或判断 $\Log$ 是否可定义。

分析方法：

\begin{itemize}
\item
  若 \(w=re^{i\theta}\ne 0\)，则
  \[
  e^z=w \Longleftrightarrow z=\log r+i(\theta+2k\pi),\quad k\in Z.
  \]
\item
  若求主值，只取 $\Arg$ 所在区间内的角。
\item
  若涉及 \(z^\alpha\)，必须先说明使用哪个对数分支。
\end{itemize}

易错点：

\begin{itemize}
\item
  \(\log z\) 不是全局单值函数。
\item
  \(\Log(zw)=\Log z+\Log w\) 不总成立，因为主值辐角可能跳过割线。
\item
  在 \(\mathbb C^*\) 上不能定义全局连续单值对数。
\end{itemize}

题型 3：判断复数列或函数项级数收敛。

分析方法：

\begin{itemize}
\item
  复数列 $z_n$ 收敛等价于 $\Re z_n$ 和 $\Im z_n$
  分别收敛。
\item
  函数项级数先找一个与点无关的上界 $M_n$，再验证
  $\sum M_n$ 收敛。
\item
  若能证明正常收敛，就自动得到局部一致收敛，后续可用于逐项求导和逐项积分。
\end{itemize}

%=====================================================================
\section{第 2 章：复微分 (Complex Differentiation)}
%=====================================================================

\subsection{主要内容}

本章的核心是复可微、解析函数、Cauchy-Riemann
方程、共形映射和调和函数。复可微的强度远高于实可微：差商极限必须与趋近方向无关，这正是
Cauchy-Riemann 方程出现的原因。

重点主线如下：

\begin{itemize}
\item
  复导数定义： \[
  f'(a)=\lim_{z\to a}\frac{f(z)-f(a)}{z-a}.
  \]
\item
  解析函数：在开集上处处复可微。
\item
  若 $f=u+iv$ 实可微，则复可微等价于 Cauchy-Riemann 方程： \[
  u_x=v_y,\qquad u_y=-v_x.
  \]
\item
  解析且 $f'(z_0)\ne 0$ 意味着局部共形。
\item
  解析函数的实部和虚部都是调和函数。
\item
  在单连通区域上，调和函数存在全局调和共轭。
\end{itemize}

\subsection{重要定理和结论}

\begin{itemize}
\item
  复可微推出连续。
\item
  复导数满足四则运算、链式法则。
\item
  常用导数： \[
  (e^z)'=e^z,\quad (\sin z)'=\cos z,\quad (\cos z)'=-\sin z,\quad (\Log z)'=\frac1z.
  \]
\item
  Cauchy-Riemann 定理：在实可微前提下，复可微等价于 \[
  u_x=v_y,\qquad u_y=-v_x.
  \]
\item
  Jacobian 的复线性形式： \[
  J_f=
  \begin{pmatrix}
  u_x&u_y\\
  v_x&v_y
  \end{pmatrix}
  =
  \begin{pmatrix}
  a&-b\\
  b&a
  \end{pmatrix}.
  \]
\item
  若 $f$ 解析，则 \[
  \det J_f=|f'(z)|^2.
  \]
\item
  共形判别：$f$ 在 $z_0$ 解析且 $f'(z_0)\ne 0$，则在
  $z_0$ 局部保持角度和方向。
\item
  若 $f=u+iv$ 解析，则 \[
  \Delta u=0,\qquad \Delta v=0.
  \]
\item
  区域上若解析函数的实部、虚部或模为常数，则函数为常数。
\end{itemize}

\subsection{典型题目分析}

题型 1：判断函数是否复可微或解析。

常见题面：判断
 \(f(z)=\bar z\)、\(f(z)=|z|^2\)、\(f(z)=x^2-y^2+2ixy\)
是否解析，并求导数。

分析方法：

\begin{itemize}
\item
  写成 $f=u+iv$。
\item
  检查 $u_x=v_y$ 和 $u_y=-v_x$。
\item
  注意 Cauchy-Riemann
  方程只在一点成立通常只能说明该点可能复可微，还需要实可微或偏导连续等条件。
\end{itemize}

典型结论：

\begin{itemize}
\item
  \(f(z)=\bar z\)：\(u=x,v=-y\)，有
  \(u_x=1\)、\(v_y=-1\)，不满足 C-R，处处不可复可微。
\item
  \(f(z)=|z|^2=x^2+y^2\)：只在
  \(0\) 处复可微，但不在任何邻域解析。
\item
  \(f(z)=z^2=x^2-y^2+2ixy\)：满足 C-R，且
  \(f'(z)=2z\)。
\end{itemize}

题型 2：求调和共轭。

常见题面：给定调和函数 $u(x,y)$，求 $v$ 使 $u+iv$
解析。

分析方法：

\begin{itemize}
\item
  使用 Cauchy-Riemann 方程： \[
  v_y=u_x,\qquad v_x=-u_y.
  \]
\item
  先对一个变量积分，再用另一条方程确定积分中的``常数函数''。
\item
  最后写出 $f=u+iv$，调和共轭只差一个实常数。
\end{itemize}

易错点：

\begin{itemize}
\item
  必须先确认定义域的拓扑条件。局部总可以找，共轭要全局存在通常需要单连通。
\item
  在 \(\mathbb C^*\)
  上，\(u=\log|z|\) 的局部共轭是
  \(\arg z\)，但不存在全局单值连续的 \(\arg z\)。
\end{itemize}

题型 3：共形映射与临界点。

常见题面：判断 $f(z)=z^n$、$e^z$、Mobius
型函数在某点是否共形，描述局部映射效果。

分析方法：

\begin{itemize}
\item
  先求 $f'(z_0)$。
\item
  若 $f'(z_0)\ne 0$，局部就是旋转加缩放：
  \[
  f(z)\approx f(z_0)+f'(z_0)(z-z_0).
  \]
\item
  若 $f'(z_0)=0$，该点是临界点，保角性失效，后续第
  4 章会用零点重数描述局部覆盖行为。
\end{itemize}

%=====================================================================
\section{第 3 章：路径积分 (Complex Integration over Paths)}
%=====================================================================

\subsection{主要内容}

本章从路径积分出发，建立 Cauchy 积分定理、Cauchy
积分公式及其推论。核心思想是：解析性可以通过积分刻画，而边界积分可以控制区域内部的函数值和导数。

重点主线如下：

\begin{itemize}
\item
  复路径积分： \[
  \int_\gamma f(z)\,dz=\int_a^b f(\gamma(t))\gamma'(t)\,dt.
  \]
\item
  有原函数、闭路积分为零、路径无关三者等价。
\item
  单连通域、星形域和多连通域上的 Cauchy 定理。
\item
  Cauchy 积分公式和广义 Cauchy 积分公式。
\item
  Cauchy 不等式、Morera 定理、Liouville 定理。
\item
  解析函数列局部一致收敛时，极限仍解析，导数也局部一致收敛。
\end{itemize}

\subsection{重要定理和结论}

\begin{itemize}
\item
  ML 估计： \[
  \left|\int_\gamma f(z)\,dz\right|\le M L(\gamma).
  \]
\item
  原函数判别：连续函数 $f$ 在区域 $D$
  上有原函数，当且仅当任意闭曲线积分为
  $0$，也等价于路径积分只依赖端点。
\item
  基本反例： \[
  \oint_{|z|=1}\frac{dz}{z}=2\pi i.
  \]
\item
  Cauchy 积分定理：在合适拓扑条件下，解析函数沿闭曲线积分为 $0$。
\item
  Cauchy 积分公式： \[
  f(z)=\frac{1}{2\pi i}\oint_{\partial U_r(a)}
  \frac{f(\zeta)}{\zeta-z}\,d\zeta.
  \]
\item
  广义 Cauchy 积分公式： \[
  f^{(n)}(z)=\frac{n!}{2\pi i}\oint_{\partial U_r(a)}
  \frac{f(\zeta)}{(\zeta-z)^{n+1}}\,d\zeta.
  \]
\item
  Cauchy 不等式： \[
  |f^{(n)}(a)|\le \frac{n!M}{r^n}.
  \]
\item
  Morera 定理：连续函数若所有三角形或可缩闭路积分为 $0$，则解析。
\item
  Liouville 定理：有界整函数必为常数。
\item
  代数基本定理可由 Liouville 定理推出。
\end{itemize}

\subsection{典型题目分析}

题型 1：直接计算路径积分。

常见题面：给定参数曲线 \(\gamma(t)\)，计算
\(\int_\gamma f(z)\,dz\)。

分析方法：

\begin{itemize}
\item
  若存在原函数 $F$，优先用 \[
  \int_\gamma f(z)\,dz=F(\gamma(b))-F(\gamma(a)).
  \]
\item
  若没有明显原函数，按定义代入参数： \[
  \int_a^b f(\gamma(t))\gamma'(t)\,dt.
  \]
\item
  若是闭路且函数在内部解析，用 Cauchy 定理直接判零。
\end{itemize}

易错点：

\begin{itemize}
\item
  $1/z$ 在不穿过 $0$ 的局部区域有原函数，但在
  \(\mathbb C^*\) 上没有全局原函数。
\item
  多连通域中，洞内奇点会贡献非零积分。
\end{itemize}

题型 2：用 Cauchy 积分公式计算围道积分。

常见题面： \[
\oint_{|z-a|=r}\frac{f(z)}{(z-b)^{n+1}}\,dz.
\]

分析方法：

\begin{itemize}
\item
  先判断 $b$ 是否在曲线内部。
\item
  若 $b$ 在内部且 $f$ 在闭圆盘邻域解析，则 \[
  \oint \frac{f(z)}{(z-b)^{n+1}}\,dz
  =\frac{2\pi i}{n!}f^{(n)}(b).
  \]
\item
  若 $b$ 在外部，且被积函数在内部解析，则积分为 $0$。
\end{itemize}

题型 3：用 Cauchy 不等式和 Liouville 定理证明函数形式。

常见题面：证明有界整函数为常数，或证明满足增长条件的整函数是多项式。

分析方法：

\begin{itemize}
\item
  对任意 $a$ 和任意大半径 $r$ 用 Cauchy 不等式。
\item
  若 \(|f(z)|\le C(1+|z|^m)\)，则对
  $n>m$ 有 \[
  |f^{(n)}(a)|\le \frac{n!M_r}{r^n}\to 0.
  \]
\item
  所以高阶导数恒为零，$f$ 是次数不超过 $m$ 的多项式。
\end{itemize}

题型 4：用 Morera 或局部一致收敛证明解析性。

分析方法：

\begin{itemize}
\item
  对极限函数，若 $f_n$ 解析且 $f_n\to f$
  局部一致，则 $f$ 解析。
\item
  对积分定义函数，若能交换极限和路径积分，则可证明三角形积分为
  $0$，再用 Morera。
\item
  题目要求``证明解析''时，不一定要直接算差商，积分刻画常更稳。
\end{itemize}

%=====================================================================
\section{第 4 章：解析函数刻画 (Characterizing Analytic Functions)}
%=====================================================================

\subsection{主要内容}

本章是整门课内容最密集的一章。它把``解析函数''从多个角度刻画：幂级数、零点、局部映射、最大模、Schwarz
引理、解析延拓、孤立奇点、Laurent 级数、留数、亚纯函数和 Mobius 变换。

重点主线如下：

\begin{itemize}
\item
  解析函数局部等于 Taylor 级数。
\item
  在开集上，复可微、Cauchy-Riemann、Cauchy 定理、Morera、Cauchy
  公式、幂级数展开等性质互相等价。
\item
  非零解析函数的零点是离散的。
\item
  解析函数由任意有内聚点的集合上的值唯一决定。
\item
  非常数解析函数是开映射，模不能在内部取最大值。
\item
  Schwarz 引理和 Schwarz-Pick 不等式控制单位圆盘自映射。
\item
  解析延拓由恒等定理保证唯一。
\item
  孤立奇点分为可去奇点、极点、本性奇点。
\item
  Laurent 级数的负幂部分决定奇点类型。
\item
  留数定理把围道积分化为奇点局部系数。
\item
  扩充复平面 \(\widehat{\mathbb C}\)
  上的亚纯函数就是有理函数。
\item
  Mobius 变换是 Riemann 球面的基本自同构。
\end{itemize}

\subsection{重要定理和结论}

\begin{itemize}
\item
  Taylor 展开： \[
  f(z)=\sum_{n=0}^{\infty}\frac{f^{(n)}(a)}{n!}(z-a)^n.
  \]
\item
  解析函数等价刻画：复可微、C-R、局部原函数、Cauchy
  积分公式、局部幂级数展开等。
\item
  初等域：任意解析函数都有全局原函数。若 $f$
  在初等域上无零点，则存在解析对数 $h$ 使 $f=e^h$。
\item
  零点分解： \[
  f(z)=(z-a)^m\varphi(z),\qquad \varphi(a)\ne 0.
  \]
\item
  恒等定理：两个解析函数若在有内聚点的集合上一致，则在整个连通区域上一致。
\item
  开映射定理：非常数解析函数把开集映为开集。
\item
  最大模原理：非常数解析函数的模不能在区域内部取最大值。
\item
  最小模原理：若非零且非常数，则模不能在内部取最小值。
\item
  Schwarz 引理：若 $f:E\to E$ 解析且
  $f(0)=0$，则 \[
  |f(z)|\le |z|,\qquad |f'(0)|\le 1.
  \]
\item
  单位圆盘自同构： \[
  \operatorname{Aut}(E)=
  \left\{e^{i\theta}\frac{z-a}{1-\bar a z}:a\in E,\theta\in R\right\}
  \] 具体符号可能因讲义中 Blaschke 因子取法而有等价变形。
\item
  Schwarz-Pick 不等式： \[
  \frac{|f'(z)|}{1-|f(z)|^2}\le \frac{1}{1-|z|^2}.
  \]
\item
  Schwarz
  反射原理：关于实轴对称的区域中，边界取实值的解析函数可通过共轭反射延拓。
\item
  Riemann 可去奇点定理：孤立奇点附近有界当且仅当可去。
\item
  极点判别：$c$ 是 $m$ 阶极点当且仅当 \[
  f(z)=\frac{h(z)}{(z-c)^m},\qquad h(c)\ne 0.
  \]
\item
  Casorati-Weierstrass 定理：本性奇点任意穿孔邻域的像在 $\C$
  中稠密。
\item
  Laurent 展开： \[
  f(z)=\sum_{j=-\infty}^{\infty}a_j(z-c)^j.
  \]
\item
  奇点分类：

  \begin{itemize}
  \item
    可去：无负幂项。
  \item
  \(m\) 阶极点：最低负幂为 \((z-c)^{-m}\)。
  \item
    本性：有无限多个非零负幂项。
  \end{itemize}
\item
  留数定义： \[
  \operatorname{Res}(f;c)=a_{-1}.
  \]
\item
  留数定理： \[
  \oint_\Gamma f(z)\,dz=2\pi i\sum_c \operatorname{Res}(f;c)
  \] 其中求和只取曲线内部奇点。
\item
  单极点留数： \[
  \operatorname{Res}(f;c)=\lim_{z\to c}(z-c)f(z).
  \]
\item
  $m$ 阶极点留数： \[
  \operatorname{Res}(f;c)=
  \frac{1}{(m-1)!}\lim_{z\to c}
  \frac{d^{m-1}}{dz^{m-1}}\left[(z-c)^m f(z)\right].
  \]
\item
  无穷远点分类：研究 $f(1/z)$ 在 $0$ 的奇点。
\item
  \(\widehat{\mathbb C}\) 上的亚纯函数当且仅当是有理函数。
\item
  $\Aut(\C)$ 为仿射映射： \[
  z\mapsto az+b,\qquad a\ne 0.
  \]
\item
  \(\operatorname{Aut}(\widehat{\mathbb C})\) 为 Mobius 变换： \[
  z\mapsto \frac{az+b}{cz+d},\qquad ad-bc\ne 0.
  \]
\end{itemize}

\subsection{典型题目分析}

题型 1：求 Taylor 展开和收敛半径。

常见题面：在某点展开
\(1/(1+z^2)\)、\(\Log(1+z)\)、\(1/(z-a)\) 等。

分析方法：

\begin{itemize}
\item
  优先化成几何级数： \[
  \frac{1}{1-w}=\sum_{n=0}^\infty w^n,\qquad |w|<1.
  \]
\item
  收敛半径由展开中心到最近奇点的距离决定。
\item
  不要只看实轴。例如 \(1/(1+x^2)\)
  在实轴无奇点，但作为复函数的奇点是 \(z=\pm i\)。
\end{itemize}

题型 2：零点、恒等定理和唯一性。

常见题面：若解析函数在一列有内聚点的点上为零或相等，证明函数恒等为零或两函数相等。

分析方法：

\begin{itemize}
\item
  先确认内聚点在定义域内部。
\item
  对差函数 $h=f-g$ 使用恒等定理。
\item
  若只是边界上有内聚点，不能直接用恒等定理，除非能先做解析延拓或反射。
\end{itemize}

典型反例：

\begin{itemize}
\item
  $\sin(\pi/z)$ 的零点 $1/n$ 聚到 $0$，但如果
  $0$ 不在定义域内，不能推出函数恒为零。
\end{itemize}

题型 3：最大模、Schwarz 引理和圆盘自同构估计。

常见题面：已知
$f:E\to E$、$f(0)=0$，证明导数或函数值估计；或已知
$f(a)=b$，求最优上界。

分析方法：

\begin{itemize}
\item
  若固定点在 $0$，直接用 Schwarz 引理。
\item
  若固定点不在 $0$，先用圆盘自同构把 $a$ 和 $b$
  移到 $0$： \[
  \phi_a(z)=\frac{z-a}{1-\bar a z}.
  \]
\item
  对 \[
  g=\phi_b\circ f\circ \phi_a^{-1}
  \] 使用 Schwarz 引理或 Schwarz-Pick。
\end{itemize}

题型 4：孤立奇点分类。

常见题面：判断
\(\sin z/z\)、\(1/(z-a)^m\)、\(e^{1/z}\)、\(\sin(1/z)\)
在某点的奇点类型。

分析方法：

\begin{itemize}
\item
  能补值且有界：可去。
\item
  趋于无穷或有限阶负幂：极点。
\item
  Laurent 展开有无限多个负幂：本性。
\end{itemize}

典型结论：

\begin{itemize}
\item
  \(\sin z/z\) 在 \(0\) 可去。
\item
  \(1/z^m\) 在 \(0\) 是 \(m\) 阶极点。
\item
  \(e^{1/z}\) 在 \(0\) 是本性奇点。
\item
  \(\sin(1/z)\) 在 \(0\) 是本性奇点，且附近零点有内聚到
  \(0\)。
\end{itemize}

题型 5：求 Laurent 展开。

常见题面：在不同环域中展开 \[
\frac{1}{z(z-1)},\qquad \frac{1}{(z-a)(z-b)}.
\]

分析方法：

\begin{itemize}
\item
  先做部分分式分解。
\item
  每一项都化成几何级数。
\item
  环域不同，展开形式不同。例如 \[
  \frac{1}{z-1}=-\frac1{1-z}=-\sum_{n=0}^\infty z^n,\quad |z|<1,
  \] 但 \[
  \frac{1}{z-1}=\frac1z\frac1{1-1/z}
  =\sum_{n=0}^{\infty}z^{-n-1},\quad |z|>1.
  \]
\end{itemize}

题型 6：用留数定理计算围道积分和实积分。

常见题面：

\begin{itemize}
\item
  计算 \(\oint f(z)\,dz\)。
\item
  计算三角有理积分 \[
  \int_0^{2\pi} R(\cos\theta,\sin\theta)\,d\theta.
  \]
\item
  计算实轴有理积分 \[
  \int_{-\infty}^{\infty}\frac{P(x)}{Q(x)}\,dx.
  \]
\item
  计算含振荡因子的积分 \[
  \int_{-\infty}^{\infty}e^{imx}\frac{P(x)}{Q(x)}\,dx.
  \]
\end{itemize}

分析方法：

\begin{itemize}
\item
  围道积分：列出曲线内部奇点，逐个算留数。
\item
  三角积分：令 \(z=e^{i\theta}\)， \[
  \cos\theta=\frac12\left(z+\frac1z\right),\quad
  \sin\theta=\frac1{2i}\left(z-\frac1z\right),\quad
  d\theta=\frac{dz}{iz}.
  \] 然后转为单位圆上的留数问题。
\item
  实轴有理积分：取上半圆或下半圆，要求大圆弧积分趋于 $0$。
\item
  含 \(e^{imz}\) 的积分：根据 \(m\)
  的符号选择上半平面或下半平面，并使用 Jordan 引理。
\end{itemize}

题型 7：无穷远点、亚纯函数和 Mobius 变换。

常见题面：判断整函数在无穷远点的奇点类型，证明
\(\widehat{\mathbb C}\) 上亚纯函数是有理函数，求 Mobius
变换或其固定点。

分析方法：

\begin{itemize}
\item
  无穷远点一律换元： \[
  \hat f(z)=f(1/z).
  \]
\item
  若整函数在无穷远点可去，则原函数有界，故为常数。
\item
  若整函数在无穷远点为极点，则原函数是多项式。
\item
  Mobius 变换由三个点的像唯一决定，求解时可用交比或直接代入线性方程。
\end{itemize}

%=====================================================================
\section{第 5 章：辐角原理与拓扑次数}
%=====================================================================

\subsection{主要内容}

本章把围道积分、零点计数和拓扑次数联系起来。核心观点是：曲线绕点的次数可以用积分表示，而解析函数的零点和极点可以通过边界像曲线的环绕数读出来。

重点主线如下：

\begin{itemize}
\item
  环绕数： \[
  w(\Gamma,a)=\frac{1}{2\pi i}\oint_\Gamma\frac{dz}{z-a}.
  \]
\item
  环绕数是整数，在 \(\mathbb C\setminus\Gamma\)
  的每个连通分支上常值。
\item
  Hopf degree theorem 描述闭曲线同伦类与环绕数的关系。
\item
  辐角原理：边界像曲线绕 $0$ 的次数等于内部零点数减极点数。
\item
  Rouche 定理用于比较两个函数的零点数。
\item
  Hurwitz 定理用于研究解析函数列极限的零点。
\item
  拓扑次数是辐角原理在更一般连续映射中的推广。
\end{itemize}

\subsection{重要定理和结论}

\begin{itemize}
\item
  环绕数整数性： \[
  w(\Gamma,a)\in Z.
  \]
\item
  环绕数的同伦不变性：曲线同伦过程中不穿过 $a$，则环绕数不变。
\item
  对正向 Jordan 曲线，内部点环绕数为 $1$，外部点为 $0$。
\item
  解析函数辐角原理： \[
  w(f\circ \Gamma,0)=N(f;\Gamma)
  \] 在无边界零点且无极点时，右边是内部零点个数，计重数。
\item
  亚纯函数辐角原理： \[
  \frac{1}{2\pi i}\oint_\Gamma \frac{f'(z)}{f(z)}\,dz
  =N(f;\Gamma)-P(f;\Gamma).
  \]
\item
  Darboux-Picard 定理：在合适条件下，边界上的单射性推出内部单射性。
\item
  Rouche 定理：若在边界
  \(|g|<|f|\)，则
  \(f\) 与 \(f+g\) 在内部零点数相同。
\item
  Hurwitz
  定理：非零解析函数列局部一致收敛时，零点在极限中稳定；若极限非零，零点不会凭空消失，只能按定理控制迁移。
\item
  拓扑次数满足规范性、可加性、同伦不变性和乘积公式。
\end{itemize}

\subsection{典型题目分析}

题型 1：计算环绕数。

常见题面：给定闭曲线 \(\Gamma\) 和点 \(a\)，求
\(w(\Gamma,a)\)。

分析方法：

\begin{itemize}
\item
  若曲线是简单正向 Jordan 曲线，直接判断 $a$ 在内部还是外部。
\item
  若曲线复杂，尝试同伦变形，只要变形中不穿过 $a$，环绕数不变。
\item
  也可以计算 \[
  \frac{1}{2\pi i}\oint_\Gamma\frac{dz}{z-a}.
  \]
\end{itemize}

易错点：

\begin{itemize}
\item
  曲线方向会改变符号。
\item
  点在曲线上时环绕数无定义。
\item
  非简单曲线的环绕数可以是任意整数，不一定只有 $0$ 或 $1$。
\end{itemize}

题型 2：用辐角原理数零点。

常见题面：求解析函数或多项式在某区域内零点个数。

分析方法：

\begin{itemize}
\item
  检查边界上 \(f(z)\ne 0\)。
\item
  计算 \[
  \frac{1}{2\pi i}\oint_\Gamma\frac{f'(z)}{f(z)}\,dz.
  \]
\item
  若直接积分困难，观察边界像曲线 \(f(\Gamma)\) 绕 \(0\)
  的次数。
\end{itemize}

典型用法：

\begin{itemize}
\item
  对多项式，边界足够大时主项控制环绕数，所以可推出代数基本定理。
\item
  对解析函数，零点必须计重数。
\end{itemize}

题型 3：用 Rouche 定理比较零点数。

常见题面：判断多项式在单位圆或某圆盘内有几个根。

分析方法：

\begin{itemize}
\item
  把函数拆成主项 $f$ 和扰动 $g$。
\item
  在边界上证明
  \(|g|<|f|\)。
\item
  得出 $f$ 和 $f+g$ 内部零点数相同。
\end{itemize}

例子模式：

\begin{itemize}
\item
  在 $|z|=1$
  上，若多项式某一项严格大于其余项模长之和，则它决定单位圆内零点数。
\item
  在大圆上，高次项通常支配；在小圆上，低次项通常支配。
\end{itemize}

题型 4：用 Hurwitz 定理处理极限函数。

常见题面：$f_n$ 解析且局部一致收敛到 $f$，已知
$f_n$ 无零点或单射，判断 $f$ 的性质。

分析方法：

\begin{itemize}
\item
  若每个 $f_n$ 无零点，则极限 $f$
  要么无零点，要么恒为零。
\item
  若每个 $f_n$ 单射，通常可推出极限 $f$
  要么单射，要么常数。
\item
  关键是局部一致收敛，不是逐点收敛。
\end{itemize}

题型 5：Darboux-Picard 与边界单射。

常见题面：给定有界域上解析函数，已知其在边界上单射，证明内部也单射。

分析方法：

\begin{itemize}
\item
  利用边界像曲线的环绕数或拓扑次数。
\item
  对任意目标值 $p$，通过边界像曲线绕 $p$
  的次数判断原像个数。
\item
  若边界单射保证边界像为 Jordan 曲线，则内部映射的重数受控。
\end{itemize}

题型 6：拓扑次数视角。

常见题面：连续映射下求方程 $h(x)=p$ 的解个数或证明解存在。

分析方法：

\begin{itemize}
\item
  若 $p$ 不在边界像上，拓扑次数 $\deg(h,\Omega,p)$
  定义良好。
\item
  同伦不穿过 $p$ 时，次数不变。
\item
  若次数非零，则方程在区域内有解。
\item
  在复分析中，解析函数的局部次数等于零点重数。
\end{itemize}

%=====================================================================
\section{复习优先级建议}
%=====================================================================

第一优先级：

\begin{itemize}
\item
  Cauchy-Riemann 方程和解析判别。
\item
  Cauchy 积分定理、Cauchy 积分公式、广义公式。
\item
  Taylor/Laurent 展开和奇点分类。
\item
  留数定理及三类积分计算。
\item
  最大模原理、恒等定理、Schwarz 引理。
\item
  辐角原理和 Rouche 定理。
\end{itemize}

第二优先级：

\begin{itemize}
\item
  解析延拓和 Schwarz 反射。
\item
  亚纯函数、无穷远点、Mobius 变换。
\item
  Hurwitz 定理、Darboux-Picard 定理。
\item
  正常收敛、Weierstrass 定理和函数列解析性。
\end{itemize}

第三优先级：

\begin{itemize}
\item
  数系构造背景。
\item
  Hopf degree theorem 和一般拓扑次数的抽象形式。
\item
  Riemann 球面的球面几何解释。
\end{itemize}

最常见的考试题链条是：

\begin{enumerate}
\def\labelenumi{\arabic{enumi}.}
\item
  先判断函数在哪些区域解析。
\item
  找奇点和边界位置。
\item
  选择 Cauchy 公式、Laurent 展开、留数定理、最大模原理或 Rouche 定理。
\item
  最后用恒等定理、Liouville 定理或辐角原理给出全局结论。
\end{enumerate}

\end{document}
